%% kapa: exam=1 round=23 year=2012 period=1 form=A subject=civil_law subject_ko=민법 seq=31 number=31
%% source: Q-Net 감정평가사 기출문제 (cst00309) artlSeq=5052578 file=raw/1008/5052578/2012년 제23회 감정평가사 제1차시험 기출문제(1교시 A형).pdf
%% conversion: kordoc 4.24.51 -> tools/gen1.py (auto); review=proofread-sonnet; flags=none
\begin{question}{31}
공유물의 분할에 관한 설명으로 옳은 것은? (다툼이 있으면 판례에 의함)
\begin{choices}
  \item 분할 협의가 성립된 경우에는 일부 공유자가 분할에 따른 이전등기에 협조하지 않더라도 공유물 분할의 소를 제기할 수 없다.
  \item 현물분할의 경우에는 공유자 전원에게 지분의 비율로 분할하고, 일부만을 공유로 남기는 방법은 허용되지 않는다.
  \item 공유물분할은 재판상 분할을 원칙으로 한다.
  \item 공유물의 분할시 각 공유자는 다른 공유자가 분할로 인하여 취득한 물건에 대하여 그 지분의 비율과 관계없이 공유물 전부에 대하여 매도인과 동일한 담보책임이 있다.
  \item 재판상 분할에 의하여 현물분할하는 경우, 분할부분의 경제적 가치의 과부족을 공유자 상호간에 금전으로 조정하게 하는 것은 허용되지 않는다.
\end{choices}
\end{question}
\answer{1}
